\section{Recorrido por el flujo de despliegue}
\label{sec:flujo_despliegue}

\subsection{Propósito del recorrido}
Este recorrido describe el flujo efectivo de despliegue para mostrar cómo se encadenan los dominios técnicos del clúster y por qué el orden evita estados intermedios inconsistentes. El foco no está en ejecutar comandos, sino en entender dependencias operativas: qué capas preparan condiciones para la siguiente, qué validaciones confirman estabilidad y dónde se aplican controles de seguridad o de idempotencia.

La secuencia se apoya en la organización real del repositorio: \texttt{site.yml} como entrada principal, \texttt{cleanup\_slurm\_gpu.yml} como flujo alternativo controlado, y roles especializados por dominio (red, almacenamiento, scheduler, GPU, entorno LLM y validación).

\subsection{Punto de entrada y criterio de orden}
El punto de entrada principal es \texttt{site.yml}. Aunque el archivo enumera etapas, el criterio útil para operación es por \textbf{macrofases}: base del sistema, conectividad y perímetro, aceleración, almacenamiento, datos de control, scheduler, entorno de aplicación y validación.

El orden está condicionado por segmentación de hosts y responsabilidades: \texttt{all} para baseline transversal, \texttt{hpc\_master} para servicios de control, \texttt{slurm\_all} para identidad/autenticación compartida, \texttt{slurm\_compute} para ejecución distribuida y \texttt{all:!hpc\_master} para comportamientos exclusivos de nodos no controladores.

\begin{filecode}{site.yml}
- name: Etapa 1 | Baseline HPC (common + ssh)
  hosts: all
  gather_facts: false
  pre_tasks:
    - ansible.builtin.setup:
        gather_subset: [all]
  roles:
    - { role: common, tags: ["common"] }
    - { role: users_ssh, tags: ["ssh"] }
\end{filecode}

\subsection{Macrofases del despliegue y dependencias}

\begin{center}
\begin{tikzpicture}[
  node distance=5mm,
  phase/.style={
    draw,
    rectangle,
    align=center,
    text width=0.70\linewidth,
    minimum height=7mm
  }
]
  \node[phase] (baseline) {Baseline};
  \node[phase, below=of baseline] (red) {Red};
  \node[phase, below=of red] (gpu) {GPU};
  \node[phase, below=of gpu] (nfs) {NFS};
  \node[phase, below=of nfs] (db) {DB};
  \node[phase, below=of db] (identidad) {Identidad};
  \node[phase, below=of identidad] (scheduler) {Scheduler};
  \node[phase, below=of scheduler] (llm) {LLM};
  \node[phase, below=of llm] (validacion) {Validación};

  \draw[-latex] (baseline) -- (red);
  \draw[-latex] (red) -- (gpu);
  \draw[-latex] (gpu) -- (nfs);
  \draw[-latex] (nfs) -- (db);
  \draw[-latex] (db) -- (identidad);
  \draw[-latex] (identidad) -- (scheduler);
  \draw[-latex] (scheduler) -- (llm);
  \draw[-latex] (llm) -- (validacion);
\end{tikzpicture}
\end{center}

\subsubsection*{1) Baseline del sistema y acceso}
Esta macrofase habilita una línea base homogénea de paquetes, utilidades y acceso remoto. Intervienen \texttt{common} y \texttt{users\_ssh} sobre \texttt{all}. Debe ejecutarse primero porque el resto del flujo asume acceso estable y servicios base disponibles. La verificación asociada se consolida en \texttt{validate}, especialmente en checks de servicio SSH y sintaxis de \texttt{sshd}.

\subsubsection*{2) Red interna y perímetro}
Aquí se establecen enlaces internos, rutas entre subredes y reglas de firewall para tráfico Slurm y administración. Intervienen \texttt{network\_internal}, \texttt{cluster\_routing} y \texttt{firewall}. Va después del baseline porque requiere interfaces y utilidades del sistema en estado consistente; y antes de Slurm/NFS porque define la conectividad que esas capas consumen. La validación se refleja en \texttt{validate/tasks/slurm.yml} mediante comprobaciones de puertos, alcance TCP y reglas esperadas.

\subsubsection*{3) Aceleración GPU/CUDA}
Esta fase instala y valida driver/toolkit NVIDIA sólo cuando el nodo realmente tiene GPU. Interviene \texttt{nvidia\_cuda}. Su ubicación en el flujo evita que capas posteriores (Slurm GRES y entorno LLM) se construyan sobre un stack de aceleración incompleto. La validación asociada aparece en \texttt{validate} (\texttt{nvidia-smi}) y en \texttt{slurm\_validate} (\texttt{srun} con \texttt{--gres}).

\subsubsection*{4) Almacenamiento compartido (NFS)}
El objetivo es exponer export NFS desde el controlador y montar en nodos cliente. Interviene \texttt{nfs\_hpc} con separación explícita por plays (server y client). Debe ocurrir antes de pruebas operativas para que rutas de trabajo compartido existan cuando las validaciones y cargas las requieran. No hay un rol de validación NFS dedicado; la estabilidad se evidencia de forma indirecta en verificaciones operativas posteriores.

\begin{filecode}{site.yml}
- name: Etapa 4 | NFS HPC (server export)
  hosts: hpc_master
  roles:
    - { role: nfs_hpc, tags: ["nfs"] }
- name: Etapa 5 | NFS HPC (clientes mount)
  hosts: "all:!hpc_master"
  roles:
    - { role: nfs_hpc, tags: ["nfs"] }
\end{filecode}

\subsubsection*{5) Base de datos y preparación de SlurmDBD}
Esta macrofase habilita persistencia de accounting. Intervienen \texttt{mariadb\_server} y luego \texttt{slurm\_db\_prep}, ambos en \texttt{hpc\_master}. El orden es estricto: primero motor y servicio, después esquema/usuarios/grants. Como validación, el rol de MariaDB impone \texttt{assert} de versión mínima y \texttt{slurm\_db\_prep} ejecuta verificaciones de base y privilegios.

\subsubsection*{6) Identidad/autenticación (Slurm identities + Munge)}
Aquí se fija consistencia de usuarios/grupos de servicio y distribución de \texttt{munge.key} en \texttt{slurm\_all}. Intervienen \texttt{slurm\_identities} y \texttt{munge}. Esta fase antecede a instalación/configuración de Slurm para evitar fallos por identidad o autenticación incompleta entre nodos. La validación se apoya en verificación local de estado Munge y en pruebas Slurm posteriores.

\subsubsection*{7) Descubrimiento de recursos (facts para Slurm)}
En \texttt{slurm\_facts} se normalizan capacidades efectivas de CPU, memoria y GPU para construir configuración declarativa coherente. Debe ejecutarse antes de \texttt{slurm\_install}/\texttt{slurm\_controller}/\texttt{slurm\_compute}, porque esos roles renderizan archivos y decisiones a partir de estos facts. La consistencia se termina de comprobar en \texttt{slurm\_validate} (particiones, estados, \texttt{srun}, smoke jobs).

\subsubsection*{8) Scheduler: control y cómputo}
Esta fase activa el plano de control y el plano de ejecución. En \texttt{hpc\_master} intervienen \texttt{slurm\_rpm\_build}, \texttt{slurm\_install} y \texttt{slurm\_controller}; en \texttt{slurm\_compute} intervienen \texttt{slurm\_install} y \texttt{slurm\_compute}. El orden evita dependencias invertidas: primero artefactos/instalación/configuración de controller, luego adopción y activación en compute. \texttt{validate} y \texttt{slurm\_validate} confirman puertos, estados de nodos, ejecución de jobs y accounting.

\subsubsection*{9) Entorno de aplicación (LLM env)}
Con \texttt{llm\_env} se instala micromamba y stack PyTorch/CUDA de forma homogénea. Esta macrofase ocurre después de CUDA y scheduler para no validar librerías sobre un runtime de GPU incompleto. Incluye validaciones internas del propio rol (por ejemplo, bloqueo de instalación de \texttt{torch} vía pip y chequeo de \texttt{torch.cuda.is\_available()}) y se complementa con \texttt{validate} cuando \texttt{validate\_llm} está habilitado.

\subsubsection*{10) Validación operativa (\texttt{validate} + \texttt{slurm\_validate})}
La validación general en \texttt{validate} corre sobre \texttt{all} y consolida salud de SSH, CUDA, Slurm y conectividad. La validación especializada en \texttt{slurm\_validate} corre en \texttt{hpc\_master} y verifica particiones, estados, \texttt{srun} CPU/GPU y smoke jobs con \texttt{sacct}. Ambas fases son de lectura y control, sin reconfiguración del estado objetivo.

\subsubsection*{11) Limpieza (playbook auxiliar) como flujo alternativo controlado}
\texttt{cleanup\_slurm\_gpu.yml} no forma parte del camino nominal de provisión; es un flujo alternativo para reprovisionar desde cero. Su diseño prioriza seguridad operativa: ejecución gradual, conservación de acceso SSH y remoción explícita de componentes Slurm/GPU/CUDA. Es útil para reiniciar estado técnico sin introducir acciones masivas simultáneas.

\subsection{Puntos críticos y decisiones de diseño}
\textbf{1) Gating temprano de topología interna con \texttt{assert}.} Antes de tocar interfaces o rutas, el flujo exige que la definición de red esté completa y tenga estructura válida; esto evita aplicar cambios parciales sobre hosts con mapeo incompleto.

\begin{filecode}{roles/network_internal/tasks/main.yml}
- name: Red Interna | Assert variables definidas
  ansible.builtin.assert:
    that:
      - network_internal_links is defined
      - network_internal_links is mapping
      - network_internal_links | length > 0
      - network_internal_keep_if is defined
      - network_internal_keep_if | length > 0
\end{filecode}

\textbf{2) Exclusión explícita de nodos sin GPU en el rol CUDA.} El despliegue evita error operacional y ruido de cambios en hosts no objetivo mediante corte de flujo por host.

\begin{filecode}{roles/nvidia_cuda/tasks/main.yml}
- name: NVIDIA/CUDA | Definir presencia de GPU
  ansible.builtin.set_fact:
    _nvidia_gpu_present: "{{ _nvidia_lspci.rc == 0 and (_nvidia_lspci.stdout | trim) != '' }}"
- name: NVIDIA/CUDA | Saltar host sin GPU
  ansible.builtin.meta: end_host
  when: not _nvidia_gpu_present
\end{filecode}

\textbf{3) Separación de NFS server/client por plays distintos.} El mismo rol \texttt{nfs\_hpc} opera en dos contextos de host diferentes; esto reduce acoplamiento y evita que lógica de export y montaje se mezcle en una sola ejecución homogénea.

\textbf{4) Aplicación de \texttt{flush\_handlers} antes de validar servicios Slurm.} Cuando hay cambios que notifican reinicios, el flujo fuerza su aplicación antes de checks de estado para no validar sobre servicios pendientes de reinicio.

\begin{filecode}{roles/slurm_compute/tasks/main.yml}
- name: Slurm | Marcar reinicio de slurmd por cambio de symlink NVML (compute)
  ansible.builtin.command: /bin/true
  changed_when: true
  notify: Slurm | Reiniciar slurmd
- name: Slurm | Aplicar reinicios pendientes antes de validar (compute)
  ansible.builtin.meta: flush_handlers
  when: slurm_manage_slurmd | bool
\end{filecode}

\textbf{5) Validaciones read-only con \texttt{changed\_when: false} y \texttt{assert}.} La verificación no introduce deriva de configuración: observa estado, evalúa criterios y falla con diagnóstico cuando se incumplen condiciones.

\begin{filecode}{roles/slurm_validate/tasks/main.yml}
- name: Validar CLI de Slurm (sinfo)
  ansible.builtin.command: sinfo --version
  register: _sinfo_ver
  changed_when: false
- name: Verificar particion CPU
  ansible.builtin.assert:
    that:
      - slurm_validate_partitions.cpu in _parts.stdout_lines
\end{filecode}

\textbf{6) Limpieza gradual con preservación de conectividad.} El flujo alternativo de limpieza usa \texttt{serial: 1} y controles explícitos para no perder acceso remoto durante remoción de servicios y reglas.

\begin{filecode}{cleanup_slurm_gpu.yml}
- name: Limpieza de Slurm y GPU/CUDA para reprovision desde cero
  hosts: all
  serial: 1
  vars:
    cleanup_preserve_ssh: true
  tasks:
    - name: Limpieza | Verificar SSH activo al finalizar
      ansible.builtin.command: systemctl is-active "{{ cleanup_ssh_service }}"
      changed_when: false
      failed_when: cleanup_ssh_active.stdout | trim != 'active'
\end{filecode}

En conjunto, estas decisiones convierten el despliegue en una secuencia controlada por dependencias reales: primero se estabilizan base y conectividad, luego se habilitan capacidades (GPU, NFS, scheduler, entorno), y al final se valida comportamiento operativo sin alterar configuración. Ese orden reduce regresiones, facilita diagnóstico y mantiene trazabilidad entre intención declarativa y estado efectivo del clúster.
